\section{Generadores tríticos, acciones cuadráticas y transporte de energía}
\label{sec:hamiltonianos-curvatura}

La geometría dinámica conserva la distinción entre el estado que genera una
realización y las coordenadas con las que ésta expresa su evolución. APP
construye y conserva las hojas aditiva y multiplicativa; el TRIT orienta el
modo operativo; el TPK compone selección, transporte y actualización sobre
el estado enriquecido mediante
\[
 U_t=\operatorname{Upd}_t\circ\operatorname{Tra}_t\circ
 \operatorname{Sel}_t.
\]
La holonomía nonádica prolonga este estado a través de
\(w_6\to w_{12}\to w_{18}\to w_{24}\to w_{30}\to R_{36}\to G_9\).
El retorno de fase conserva el avance de memoria. Sus representaciones
\(\pi_{\rm HMT},\varphi_{\rm HMT},e_{\rm HMT},\alpha_{\rm HMT}\)
pertenecen a la misma generación correlacionada, con los lectores y el
orden constructivo interno ya establecidos. La estructura discreta conjunta
del continuo sostiene después las representaciones geométricas y físicas.
En este punto de la construcción, un Hamiltoniano es el funcional que
realiza un transporte ya tipado sobre una fibra simpléctica; su dominio,
su orientación y su parámetro temporal forman parte de la afirmación.

El artículo sobre la extrema y media razón holográfica recupera tres
operadores concretos. La multiplicación por \(4\) módulo \(9\) actúa sobre
las órbitas ternarias de unidades; en su módulo de aumento tiene matriz
\[
 C=\begin{pmatrix}0&-1\\1&-1\end{pmatrix},\qquad C^2+C+I=0.
\]
El transporte parabólico y la incidencia areal--volumétrica están
representados, en sus respectivos planos, por
\[
 N=\begin{pmatrix}0&1\\0&0\end{pmatrix},\qquad
 F_{\rm av}=\begin{pmatrix}0&1\\1&1\end{pmatrix}.
\]
La incidencia procede del calendario de hojas; sus potencias cuentan
composiciones de esos modos. Los generadores reales expresados son
\begin{equation}
 J_+=\frac{2C+I}{\sqrt3},\qquad J_0=N,\qquad
 J_-=\frac{2F_{\rm av}^{2}-3I}{\sqrt5},\qquad
 J_\tau^2=-\tau I.
 \label{eq:hc-generadores}
\end{equation}
Aquí \(\tau=+1,0,-1\) discrimina el régimen elíptico, parabólico e
hiperbólico. Las fibras tienen procedencias y bases propias. Las fórmulas
\(e^{\pi J_+}=-I\), \(e^{J_0}=I+J_0\) y
\(e^{2\log\varphi J_-}=F_{\rm av}^{2}\) evalúan coordenadas HMT
previamente generadas en esos operadores~\cite{hmtX}.

\subsection{Forma alternante y Hamiltoniano del generador}

En una fibra real orientada de dimensión dos se fija
\[
 \omega=dQ\wedge dP,\qquad
 \Omega=\begin{pmatrix}0&1\\-1&0\end{pmatrix},\qquad z=(Q,P)^T.
\]
Se adopta \(\iota_{X_H}\omega=dH\); por tanto
\(\dot Q=\partial_PH\), \(\dot P=-\partial_QH\).
Esta convención fija también el signo de la acción orientada.

\begin{proposition}[Realización hamiltoniana de un generador de traza nula]
\label{prop:hc-cuadratica}
Para una matriz real \(J\) de tamaño dos y traza nula, la matriz
\(G_J=-\Omega J\) es simétrica y el Hamiltoniano
\begin{equation}
 H_J(z)=\frac12z^TG_Jz
 \label{eq:hc-hj}
\end{equation}
produce exactamente \(\dot z=Jz\). Si \(J^2=-\tau I\),
entonces \(\det G_J=\tau\). Para un isomorfismo simpléctico
\(B:V\to V'\), el transporte
\[
 J'=BJB^{-1},\qquad H'(z')=H_J(B^{-1}z')
\]
satisface \(G'=B^{-T}G_JB^{-1}\) y entrelaza ambos flujos.
\end{proposition}
\begin{proof}
Escribiendo \(J=\bigl(\begin{smallmatrix}a&b\\c&-a\end{smallmatrix}\bigr)\),
se obtiene \(G_J=\bigl(\begin{smallmatrix}-c&a\\a&b\end{smallmatrix}\bigr)\).
Así \(\Omega\nabla H_J=\Omega(-\Omega J)z=Jz\).
Cayley--Hamilton da \(J^2=-\det(J)I\), y
\(\det(-\Omega)=1\). Finalmente, \(B^T\Omega B=\Omega\)
permite sustituir \(z=B^{-1}z'\) tanto en la forma como en la ecuación.
\end{proof}

Aplicada a \eqref{eq:hc-generadores}, la proposición da
\begin{equation}
 H_{J_+}=-\frac{Q^2-QP+P^2}{\sqrt3},\qquad
 H_{J_0}=\frac{P^2}{2},\qquad
 H_{J_-}=\frac{-Q^2-QP+P^2}{\sqrt5}.
 \label{eq:hc-hamiltonianos-nativos}
\end{equation}
El signo elíptico registra la orientación de \(C\) en la base de aumento.
El flujo inverso \(-J_+\) posee la forma positiva opuesta. La elipse
de acción del corpus utiliza precisamente una orientación de fase
decreciente para su Hamiltoniano positivo. Conservar este signo hace
compatible el operador cíclico con la acción \(P\,dQ\).

La conjugación también proporciona un criterio exacto de comparación.
Si \(BJ_\tau=J_{\tau'}B\) y \(B\) es invertible, elevar al cuadrado
da \((\tau-\tau')B=0\), luego \(\tau=\tau'\). Un enlace puede
conservar \(\omega\) entre fibras de regímenes distintos y conservar
sus memorias respectivas; la conjugación de los generadores exige además
igualdad del régimen. Esta diferencia permite estudiar una transición
trítica sin sustituirla por un cambio de nombre del mismo flujo.

\subsection{Deformación de la elipse y transformación de Legendre}

La coordenada \(\eta=\operatorname{artanh}(C^*/A)\), en el dominio
\(|C^*/A|<1\), define el marco simpléctico
\(M_\eta=\operatorname{diag}(e^{\eta/2},e^{-\eta/2})\).
La deformación de las dos direcciones es recíproca y conserva su producto.
En la coordenatización normalizada \((q,p)=M_\eta^{-1}(Q,P)\), consideremos
\[
 L_\tau=\begin{pmatrix}0&1\\-\tau&0\end{pmatrix},\qquad
 I_{\tau,\eta}=\frac12\bigl(\tau e^{-\eta}Q^2+e^\eta P^2\bigr).
\]
Estos representantes clasifican flujos cuadráticos orientados. Su empleo
en una fibra nativa se hace mediante un marco simpléctico declarado,
conforme a la proposición~\ref{prop:hc-cuadratica}. El representante
elíptico positivo es el usado por la elipse de acción; para él, \(I\)
tiene dimensión de acción.

Los cambios de marco pueden escribirse explícitamente. Si
\(b_+=\sqrt{\sqrt3/2}\) y \(b_-=\sqrt{\sqrt5/2}\), las matrices
\[
 B_+=\begin{pmatrix}b_+&b_+/\sqrt3\\0&2b_+/\sqrt3\end{pmatrix},
 \qquad
 B_-=\begin{pmatrix}b_-&-b_-/\sqrt5\\0&2b_-/\sqrt5\end{pmatrix}
\]
tienen determinante uno y satisfacen
\( (-J_+)B_+=B_+L_{+1}\) y
\(J_-B_-=B_-L_{-1}\). Para el régimen parabólico se toma
\(B_0=I\). La multiplicación de las dos columnas verifica las
igualdades, y el determinante uno prueba la conservación de la forma
alternante.

\begin{theorem}[Transporte variable y lagrangiano regular]
\label{thm:hc-legendre}
Sean \(\eta(t)\) de clase \(C^1\) y \(\omega(t)>0\). El Hamiltoniano
\begin{equation}
 H_{\tau,\rm tr}
 =\omega I_{\tau,\eta}+\frac{\dot\eta}{2}QP
 \label{eq:hc-htr}
\end{equation}
conserva \(I_{\tau,\eta(t)}\) a lo largo de sus soluciones. Su
transformación de Legendre respecto de \(P\) es invertible y da
\begin{equation}
 \mathcal L_{\tau,\rm tr}(Q,\dot Q,t)
 =\frac{(\dot Q-\dot\eta Q/2)^2}{2\omega e^\eta}
 -\frac{\omega\tau e^{-\eta}}2Q^2.
 \label{eq:hc-ltr}
\end{equation}
\end{theorem}
\begin{proof}
Las ecuaciones son
\[
 \dot Q=\omega e^\eta P+\frac{\dot\eta}2Q,
 \qquad
 \dot P=-\omega\tau e^{-\eta}Q-\frac{\dot\eta}2P.
\]
Derivar \(I_{\tau,\eta(t)}\) y sustituirlas cancela por separado
los términos \(\omega QP\) y los términos de variación de \(\eta\).
La derivada segunda \(\partial_P^2H=\omega e^\eta>0\)
permite resolver
\(P=(\dot Q-\dot\eta Q/2)/(\omega e^\eta)\).
Sustituir en \(P\dot Q-H\) prueba \eqref{eq:hc-ltr}.
Además,
\[
 P\dot Q-H_{\tau,\rm tr}
 =p\dot q-\frac\omega2(\tau q^2+p^2),
\]
lo que prueba que el término mixto es la conexión variacional del marco
móvil. La conservación procede del transporte completo y no de mantener
constantes los semiejes.
\end{proof}

El resultado elíptico y su acción están desarrollados en el corpus
variacional~\cite{hmtX}; la misma prueba exhibe sus dos realizaciones
cuadráticas conjugadas. La nulidad del determinante de la forma parabólica
convive con una transformación de Legendre regular: ésta depende de la
derivada segunda en el momento, no del determinante de la forma completa.
En cambio, para el levantamiento cotangente lineal
\(H_A(q,p)=p^TAq\), se tiene \(\partial_p^2H_A=0\).
Su acción de primer orden conserva \(\dot q=Aq\) mediante un
multiplicador; es un objeto variacional distinto de
\eqref{eq:hc-ltr}.

\subsection{Retorno de memoria y carga variacional conservada}

En el registro de memoria del calendario, los contadores \((g,s)\)
reciben el incremento \((4,72)\) en el retorno completo. La relación
\(72=18\cdot4\) fija
\[
 v=(1,18)^T,\qquad I_{\rm mem}=s-18g,
 \qquad B=\begin{pmatrix}1&0\\18&1\end{pmatrix}.
\]
El estabilizador nativo es
\begin{equation}
 N_v=BNB^{-1}=
 \begin{pmatrix}-18&1\\-324&18\end{pmatrix},\qquad
 N_vm=vI_{\rm mem}(m),\qquad N_v^2=0.
 \label{eq:hc-nv}
\end{equation}
En el plano realificado \(m=(g,s)^T\), dotado de \(dg\wedge ds\),
el Hamiltoniano \(H_{\rm mem}=I_{\rm mem}^2/2\) produce
\(dm/du=N_vm\). El parámetro \(u\) pertenece a esta realización
adimensional del registro. La coordenatización
\(q=g,\ p=s-18g\) es simpléctica y transforma la uno-forma de acción
según
\begin{equation}
 s\,dg=p\,dq+d(9q^2).
 \label{eq:hc-memoria-borde}
\end{equation}
Por tanto, la eliminación de \(s\) en la acción de primer orden da
\begin{equation}
 \mathcal L_{\rm mem}
 =\frac12\left(\frac{dg}{du}\right)^2
 +18g\frac{dg}{du}
 =\frac12\left(\frac{dg}{du}\right)^2
 +\frac{d}{du}(9g^2).
 \label{eq:hc-lmem}
\end{equation}
La ecuación variacional es \(d^2g/du^2=0\), y el momento
conservado en la coordenatización simpléctica es \(p=I_{\rm mem}\).
El retorno afín completo procede del Hamiltoniano lineal
\(H_{\rm ret}=4I_{\rm mem}\): su aplicación al tiempo \(u=1\)
es \((g,s)\mapsto(g+4,s+72)\). Ambos Hamiltonianos conmutan
porque son funciones de la misma carga, pero realizan operaciones
distintas: estabilización parabólica y traslación del registro.
Estas fórmulas recuperan el mecanismo variacional de la memoria, incluido
su término de frontera~\cite{hmtX}.

\subsection{Curvatura métrica y energía geodésica}

La realización métrica de los regímenes utiliza una coordenatización radial y una
escala inversa de longitud \(\rho>0\):
\begin{equation}
 g_{\tau,\rho}=dr^2+S_{\tau,\rho}(r)^2d\theta^2,
 \qquad
 S_{\tau,\rho}(r)=
 \begin{cases}
 \sin(\rho r)/\rho,&\tau=+1,\\
 r,&\tau=0,\\
 \sinh(\rho r)/\rho,&\tau=-1.
 \end{cases}
 \label{eq:hc-metrique}
\end{equation}
La curvatura gaussiana es \(K_G=-S''/S=\tau\rho^2\).
La coordenatización se considera en su interior regular, \(S>0\); para la
realización esférica, \(0<r<\pi/\rho\). Los polos se tratan con
coordenatizaciones regulares diferentes. El parámetro \(\rho\) y la realización
dimensional de la métrica mantienen su declaración propia.

Para una masa realizada \(m>0\), el lagrangiano geodésico y su
Hamiltoniano son
\begin{equation}
 \mathcal L_{\rm geo}=\frac m2(\dot r^2+S^2\dot\theta^2),\qquad
 H_{\rm geo}=\frac1{2m}\left(p_r^2+\frac{p_\theta^2}{S^2}\right).
 \label{eq:hc-geodesique}
\end{equation}
En efecto, \(p_r=m\dot r\), \(p_\theta=mS^2\dot\theta\);
el hessiano de velocidades es \(m\operatorname{diag}(1,S^2)\)
y es positivo e invertible. La ecuación radial es
\(\ddot r-SS'\dot\theta^2=0\), mientras
\(d(mS^2\dot\theta)/dt=0\) conserva el momento angular.
La positividad de esta energía vale para las tres curvaturas. La forma
indefinida del generador hiperbólico de \eqref{eq:hc-hamiltonianos-nativos}
y la energía geodésica positiva de una métrica de curvatura negativa
son funciones sobre espacios distintos. Esta comparación fija exactamente
qué significado de curvatura interviene en cada afirmación.

\subsection{Respuesta constitutiva y realización modal de Maxwell}

Las secciones de acción y las coordenadas angulares preceden a la
respuesta constitutiva \cite{hmtIntegral}. En la orientación real fijada \(x>y>0\),
con los proyectores de las dos hojas,
\[
 T=q_+P_++q_-P_-,\quad q_\pm=e^{-x\mp y},\quad
 \mathcal V=(I-T^{90})(I-T^{120})^{-1},
 \qquad 0<q_+<q_-<1,
\]
el cálculo funcional produce
\[
 r_\pm=\frac{1-q_\pm^{90}}{1-q_\pm^{120}},\qquad
 \widehat\varepsilon=r_+^2,\quad \widehat\mu=r_-^2,
 \quad \widehat c=(r_+r_-)^{-1},\quad \widehat Z=r_-/r_+.
\]
Las coordenatizaciones dimensionales correspondientes conservan exactamente
\(\varepsilon\mu=c^{-2}\) y \(\mu/\varepsilon=Z^2\).
La permitividad, la permeabilidad y la velocidad son así representaciones
del mismo operador, con los sectores de incidencia fijados por el TPK;
la dinámica siguiente las emplea como resultados ya construidos.

Para exhibir el transporte hacia un Hamiltoniano de campo, se fija una
realización espacial orientada, condiciones de frontera que permitan
integración por partes, y un modo transversal real \(f\) normalizado por
\[
 \nabla\cdot f=0,\qquad \int|f|^2=1,\qquad
 \nabla\times\nabla\times f=k^2f,\qquad k>0.
\]
En calibre transversal y sin fuentes, \(A=af\) y el momento
\(\Pi=pf\) definen una inmersión del plano modal en el espacio de
campos. Para constantes constitutivas homogéneas, el Hamiltoniano y la
forma simpléctica de Maxwell se restringen a
\begin{equation}
 H_{\rm EM}(a,p)=\frac{p^2}{2\varepsilon}
 +\frac{k^2a^2}{2\mu},\qquad \omega_{\rm EM}=da\wedge dp.
 \label{eq:hc-maxwell-mode}
\end{equation}
La integración por partes da
\(\int|\nabla\times f|^2=k^2\), que prueba esta restricción.
Con \(\omega_k=ck\), el cambio
\begin{equation}
 Q=\sqrt{\varepsilon\omega_k}\,a,
 \qquad P=\frac{p}{\sqrt{\varepsilon\omega_k}}
 \label{eq:hc-maxwell-symplectic}
\end{equation}
preserva \(da\wedge dp=dQ\wedge dP\) y da
\(H_{\rm EM}=\omega_k(Q^2+P^2)/2\). Su inversa y sus ecuaciones
\(\dot Q=\omega_kP,\ \dot P=-\omega_kQ\)
constituyen el entrelazamiento modal con el régimen elíptico positivo.
La restricción es invariante porque el operador de Maxwell conserva el
modo propio transversal. La transformación de Legendre produce
\(\mathcal L_{\rm EM}=\varepsilon\dot a^2/2-k^2a^2/(2\mu)\).

Dos modos transversales de igual frecuencia, combinados en cuadratura,
realizan las dos orientaciones circulares del corpus. El campo satisface
\(\mathbf B=c^{-1}\widehat{\mathbf k}\times\mathbf E\) y
\(\omega=ck\). La relación con el oscilador se demuestra aquí mediante
la inmersión de campos, la normalización modal y el transporte de la forma
simpléctica, manteniendo explícitas las condiciones de frontera.
Un cambio de \(k\), de la geometría de la cavidad o de la preparación
modifica el modo realizado; mantiene intacta la genealogía de
\(\varepsilon,\mu,c,Z\).

En la coordenatización lorentziana del mismo sector, el acoplamiento a una sección
material se escribe con \(F=dA\) y una carga \(q_e\) ya tipada.
Para esta realización variacional se toma \(\nabla\) igual a la
conexión de Levi--Civita de \(g\). Sobre
\(T^*\mathcal M\), el Hamiltoniano covariante
\[
 H_A=\frac1{2m}g^{\mu\nu}(p_\mu-q_eA_\mu)(p_\nu-q_eA_\nu)
\]
es la transformación de Legendre de
\(\mathcal L_A=\frac m2g_{\mu\nu}\dot x^\mu\dot x^\nu
+q_eA_\mu\dot x^\mu\). El hessiano \(mg\) es invertible.
La ecuación de Euler--Lagrange da
\(m\nabla_u u^\mu=q_eF^\mu{}_{\nu}u^\nu\).
La antisimetría de \(F\) conserva \(g(u,u)\). Se trata de la
realización covariante con parámetro afín, mientras
\eqref{eq:hc-maxwell-mode} utiliza el tiempo de una descomposición
espacial transversal: ambos parámetros y ambos dominios quedan definidos.

\subsection{Canales graduados, representación de Clifford y operador de masa}

El atlas material produce primero canales completos de ruta, hoja, sabor,
color y representación. En una profundidad finita, su realización
hermítica es \(\mathcal H_K=\bigoplus_{s\in S_K}V_s\), con base
ortonormal de canales y graduación de paridad transportada. El álgebra
exterior del sector impar genera las CAR: la inserción de un modo
\(a_j^*\) y su retirada adjunta \(a_j\) satisfacen
\(\{a_i,a_j^*\}=\delta_{ij}I\) y \(\{a_i,a_j\}=0\).
La prueba cuenta los signos de intercambio de modos completos; las masas
se evalúan después~\cite{hmtVI}.

En el sector exterior de un solo modo, la base \((|0\rangle,|1\rangle)\)
identifica el operador de retirada con
\(a=N\). El mapa \(e_1\mapsto|0\rangle,\ e_2\mapsto|1\rangle\)
es una isometría del plano complejo normalizado y entrelaza ambos
nilpotentes. Esta identificación emplea el producto interno y la
graduación declarados; conserva el distinto significado de sus bases.
Con el \(i\) interno ya realizado, se obtienen
\begin{equation}
 \sigma_1=N+N^*,\qquad
 \sigma_2=-i(N-N^*),\qquad
 \sigma_3=NN^*-N^*N.
 \label{eq:hc-pauli}
\end{equation}
La multiplicación directa da
\(\sigma_j^*=\sigma_j\) y
\(\sigma_i\sigma_j=\delta_{ij}I+i\epsilon_{ijk}\sigma_k\).
En particular, la pareja nilpotente y su adjunto producen una
representación de Clifford; el adjunto es el de la fibra hermítica y
se transporta junto con ella.

El operador de masa sobre una colección finita de canales \(\mathcal F\)
está construido por el registro, el carácter y sus refinamientos:
\begin{equation}
 M_\beta=S M_0 S,\qquad
 S=R_{\rm act}^{B_M/2},\quad
 R_{\rm act}=\hbar_{\rm pre}/\hbar_{\rm ret}>0,
 \qquad M_0>0.
 \label{eq:hc-mass}
\end{equation}
Aquí \(B_M\) es el operador diagonal del exponente de ruta y \(M_0\)
contiene el carácter refinado relativo al origen electrónico.
La positividad sigue de \(\langle v,M_\beta v\rangle
=\langle Sv,M_0Sv\rangle>0\) para \(v\ne0\).
El producto de hojas conserva el sector nulo por su realización propia;
la exponencial del carácter material permanece positiva.

En esta realización se escribe \(\hbar:=\hbar_{\rm ret}\),
la sección de acción ya construida, y \(c:=c_{\rm vac}\) en la
coordenatización dimensional declarada. Sobre \(\mathbb C^4\otimes\mathcal F\), definamos
\[
 \alpha_j^{D}=\sigma_1\otimes\sigma_j,\qquad
 \beta^D=\sigma_3\otimes I_2.
\]
El superíndice distingue estas matrices de la constante de estructura fina.
La coordenatización espacial \(1+3\) asigna tres componentes de momento
\(\mathbf p\). La aplicación de multiplicación
\begin{equation}
 H_D(\mathbf p)
 =c\sum_{j=1}^3\alpha_j^Dp_j\otimes I_{\mathcal F}
 +c^2\beta^D\otimes M_\beta
 \label{eq:hc-dirac}
\end{equation}
es la realización libre de Dirac de ese operador material.

\begin{theorem}[Factorización espinorial de la dispersión material]
\label{thm:hc-dirac-square}
La matriz \eqref{eq:hc-dirac} es autoadjunta y satisface
\begin{equation}
 H_D(\mathbf p)^2
 =I_4\otimes\bigl(c^2|\mathbf p|^2I_{\mathcal F}
 +c^4M_\beta^2\bigr).
 \label{eq:hc-dirac-square}
\end{equation}
En un autoespacio \(M_\beta v=m v\), sus energías son
\(\pm E_m(\mathbf p)\), donde
\(E_m=\sqrt{c^2|\mathbf p|^2+m^2c^4}\).
\end{theorem}
\begin{proof}
Las relaciones de \eqref{eq:hc-pauli} dan
\(\{\alpha_i^D,\alpha_j^D\}=2\delta_{ij}I_4\),
\(\{\alpha_i^D,\beta^D\}=0\) y \((\beta^D)^2=I_4\).
El factor de sabor conmuta con las matrices espinoriales.
Los términos cruzados se anulan y los cuadrados dan
\eqref{eq:hc-dirac-square}. La traza nula y las relaciones de Clifford
producen ambas ramas, con multiplicidad espinorial dos cuando \(E_m>0\).
\end{proof}

En representación de momentos, el dominio natural es
\(\{\psi\in L^2(\mathbb R^3;\mathbb C^4\otimes\mathcal F):
H_D\psi\in L^2\}\). Al ser una matriz hermítica mensurable actuando
por multiplicación, el operador es autoadjunto en ese dominio. La
representación espacial procede de la transformación unitaria de Fourier
posterior a la construcción de los canales. El teorema fija la realización
libre y su dependencia exacta de la masa HMT; la interacción conserva sus
operadores de conexión y su dominio propio.

En la rama positiva \(m>0\), la transformación de Legendre ordinaria
del símbolo \(E_m(\mathbf p)\) es
\[
 \mathbf v=\frac{c^2\mathbf p}{E_m},\qquad
 \mathcal L_m=-mc^2\sqrt{1-|\mathbf v|^2/c^2},
 \qquad |\mathbf v|<c.
\]
Resolver \(\mathbf p=m\mathbf v/\sqrt{1-|\mathbf v|^2/c^2}\)
y sustituir en \(\mathbf p\cdot\mathbf v-E_m\) demuestra la
igualdad. La acción espinorial, en cambio, es de primer orden en el campo:
\(\mathcal L_D=\frac{i\hbar}{2}(\psi^*\dot\psi-\dot\psi^*\psi)
-\psi^*H_D\psi\). Su variación da
\(i\hbar\dot\psi=H_D\psi\); su hessiano en velocidades del campo
es singular. Ambas afirmaciones son compatibles porque actúan sobre
espacios variacionales diferentes.

Para una conexión abeliana estática en una coordenatización espacial plana, con
\(\Pi_j=-i\hbar\partial_j-q_eA_j\), masa constante y potencial
escalar nulo, se conserva sobre funciones lisas de soporte compacto
\[
 [\Pi_i,\Pi_j]=iq_e\hbar\epsilon_{ijk}B_k,
\]
y la misma factorización produce
\begin{equation}
 H_D(A)^2=c^2\boldsymbol\Pi^2+c^4M_\beta^2
 -q_e\hbar c^2\boldsymbol\Sigma\cdot\mathbf B,
 \qquad \Sigma_j=I_2\otimes\sigma_j.
 \label{eq:hc-dirac-curvature}
\end{equation}
Los factores identidad de sabor quedan implícitos. El término de curvatura
es el producto de los dos antisimétricos: el conmutador de derivadas
covariantes y la multiplicación de Pauli. Su coeficiente queda fijado
al elegir la representación de carga y la coordenatización de acción; no se
selecciona mediante un valor espectral objetivo.

Si \(V\) es la transformación unitaria de sabor ya construida, entonces
\[
 (I_4\otimes V)H_D(M_\beta)(I_4\otimes V^*)
 =H_D(VM_\beta V^*).
\]
La mezcla unitaria conserva el espectro completo. La congruencia positiva
de \eqref{eq:hc-mass} transporta otra estructura y puede modificar los
autovalores en una métrica fija. Por ejemplo,
\(M_0=\operatorname{diag}(1,2)\), \(S=\operatorname{diag}(2,1)\)
da \(SM_0S=\operatorname{diag}(4,2)\). Si se transporta también el
producto interno a \(G'=S^*S\), el problema generalizado
\(SM_0Sv=\lambda G'v\) recupera los autovalores de \(M_0\)
mediante \(w=Sv\). De este modo, cambio de base de sabor, deformación
material y transporte de métrica tienen consecuencias espectrales precisas.

\subsection{Curvatura del potencial escalar en la realización electrodébil}

El sector electrodébil utiliza el determinante del mismo operador angular,
\(D_A=\det T=e^{-2A_{\rm rad}}\), junto con las firmas materiales ya
construidas. En la realización angular con refinamientos comunes de
\(W\) y \(Z\), sus firmas \((94,-1)\) y \((95,-1)\) dan
\(m_W^\angle/m_Z^\angle=e^{-A_{\rm rad}}\).
En la convención racionalizada de árbol del corpus,
\[
 g^2=\frac{4\pi_{\rm HMT}\alpha_{\rm HMT}}{1-D_A},\qquad
 \lambda_H=\frac{\pi_{\rm HMT}\alpha_{\rm HMT}}{2(1-D_A)}
 \exp\!\left(6A_{\rm rad}+\frac{10}3C^*_{\rm rad}\right)>0.
\]
La segunda fórmula resulta de la firma escalar \((97,9)\) y de
\(m_H^2/m_W^2=8\lambda_H/g^2\). La normalización del potencial es
\(V(\mathsf H)=-\mu_H^2\mathsf H^*\mathsf H
+\lambda_H(\mathsf H^*\mathsf H)^2\). Estas composiciones conservan
su escala y su convención de árbol; los refinamientos que se cancelan
en una razón quedan declarados en la propia realización.

En una coordenatización radial real \(\mathsf H^*\mathsf H=h^2/2\), el potencial
es \(V(h)=-\mu_H^2h^2/2+\lambda_Hh^4/4\). Para
\(\mu_H^2>0\), sus puntos críticos son \(0\) y
\(\pm v\), con \(v^2=\mu_H^2/\lambda_H\), y
\[
 V''(0)=-\mu_H^2,\qquad V''(\pm v)=2\mu_H^2.
\]
En unidades naturales de esta coordenatización, el modo homogéneo normalizado posee
\(H_h=p_h^2/2+V(h)\). Su linealización en un punto crítico \(h_0\)
tiene generador
\[
 A_{h_0}=\begin{pmatrix}0&1\\-V''(h_0)&0\end{pmatrix},
 \qquad A_{h_0}^2=-V''(h_0)I.
\]
Por tanto, el origen es hiperbólico y los mínimos son elípticos; la
frontera \(V''=0\) tiene el tipo parabólico. Esta afirmación se deduce
de la derivada segunda del potencial y de una normalización modal
explícita. La energía material positiva del mínimo y la dirección
inestable de otro fondo pertenecen así a la misma función con distintos
puntos de expansión. El modo linealizado es un transporte particular de
los tres tipos cuadráticos; sus frecuencias y unidades proceden del
potencial, no de identificar su hessiano con la curvatura gaussiana de
\eqref{eq:hc-metrique}.

\subsection{Modo de curvatura de Yang--Mills y oscilación elíptica}

La realización de Yang--Mills construida en las secciones precedentes
posee un generador positivo \(D^{\rm YM}\) de dimensión inversa de
longitud. La escala
\[
 \kappa_{\rm av}=\frac{\mu_{\rm vol}}{\ell_0}
 \left(\frac{\pi_{\rm HMT}}{\varphi_{\rm HMT}^2}-1\right)>0,
 \qquad \delta=3\kappa_{\rm av}
\]
desciende de la incidencia de hojas y del lector regional, con la unidad
longitudinal \(\ell_0\) declarada. El vacío es simple y el testigo de
curvatura en la coordenatización producto es
\[
 W(q)=\mathbf1_{\{q_1+\cdots+q_6=0\}}-\frac13,
 \quad \mathbb EW=0,\quad \|W\|^2=\frac29,
 \quad D^{\rm YM}W=\delta W.
\]
El transporte unitario hacia la medida de interacción lleva \(W\) a
\(W_g\) y conserva esas identidades. Las inclusiones de refinamiento
\(J_m\) entrelazan los generadores y conservan el testigo, de modo que
éste permanece en el límite físico. Aquí se utiliza el resultado completo
de las pruebas de vacío, brecha uniforme y reconstrucción, no un
autovalor aislado de una matriz finita~\cite{HMT-YM}.

Sea \(w=W_g/\sqrt{2/9}\) el vector unitario en la fibra física y
\(\hbar=\hbar_{\rm ret}>0\). El Hamiltoniano energético es
\(H_E=\hbar cD^{\rm YM}\). Sobre el espacio de Hilbert realificado
se adopta la forma \(\omega_{\mathcal H}(u,v)=2\hbar
\operatorname{Im}\langle u,v\rangle\), con producto interno antilineal
en el primer argumento. Definamos
\begin{equation}
 \mathfrak F(Q,P)=\frac{Q+iP}{\sqrt{2\hbar}}\,w.
 \label{eq:hc-ym-embedding}
\end{equation}

\begin{theorem}[Inmersión simpléctica del modo físico de curvatura]
\label{thm:hc-ym-oscillator}
La aplicación \(\mathfrak F:\mathbb R^2\to\mathcal H^{\rm phys}\)
es una inmersión real con imagen invariante por la evolución física y
satisface
\begin{equation}
 \mathfrak F^*\omega_{\mathcal H}=dQ\wedge dP,
 \qquad
 \langle\mathfrak F,H_E\mathfrak F\rangle
 =\frac{c\delta}{2}(Q^2+P^2).
 \label{eq:hc-ym-energy-pullback}
\end{equation}
La ecuación de Schrödinger restringida entrelaza exactamente el flujo
elíptico \(\dot Q=c\delta P,\ \dot P=-c\delta Q\).
El mapa es compatible con las inclusiones de refinamiento.
\end{theorem}
\begin{proof}
Para dos variaciones \((a,b),(a',b')\), el producto interno de sus
imágenes tiene parte imaginaria \((ab'-ba')/(2\hbar)\).
Esto prueba la primera identidad. La ecuación propia y \(\|w\|=1\)
dan la segunda. Puesto que
\(\dot\psi=-icD^{\rm YM}\psi\), el coeficiente \(Q+iP\)
satisface \(\dot Q+i\dot P=-ic\delta(Q+iP)\).
Finalmente, \(J_mw_m=w_{m+1}\) implica
\(J_m\mathfrak F_m=\mathfrak F_{m+1}\). Estas identidades pasan al
límite por las isometrías ya construidas.
\end{proof}

El plano completo describe amplitudes de un modo; su círculo
\(Q^2+P^2=2\hbar\) representa los vectores normalizados de ese
autoespacio complejo. Sobre tal círculo la fase global evoluciona y el
rayo cuántico permanece fijo. Superposiciones de autoespacios distintos
conservan, además, sus fases relativas. La inmersión mantiene esta
diferencia entre espacio de amplitudes y espacio proyectivo.

La energía asociada al salto espectral y el salto de masa son
\(\Delta_E=\hbar c\delta\) y \(m_{\rm gap}=\hbar\delta/c\).
La frecuencia modal \(c\delta\) tiene dimensión de tiempo inverso.
Las longitudes euclídeas de \(e^{-sD^{\rm YM}}\) y el tiempo físico
de \(e^{-itH_E/\hbar}\) se relacionan mediante la coordenatización y la
reconstrucción, conservando sus distintos grupos y semigrupos.
La composición \(\mathfrak F\circ M_\eta^{-1}\) transporta esta
misma energía a la elipse de acción; si el marco varía, el término de
conexión de \eqref{eq:hc-htr} completa su evolución.

El enlace obtenido es una realización exacta sobre el modo invariante de
curvatura y sus prolongaciones. La positividad del Hamiltoniano cuántico
genera una rotación en su realificación simpléctica; la denominación
hiperbólica de una métrica espacial pertenece al tensor métrico, y el
transporte unipotente del registro pertenece a su memoria. En dimensión
finita, un grupo \(I+tN\) unitario para todo \(t\in\mathbb R\), con
\(N^2=0\), fuerza \(N=0\): derivar la unitariedad da \(N^*=-N\),
y entonces \(N^*N=-N^2=0\). El flujo parabólico clásico conserva su
forma simpléctica y posee su cuantización en espacios apropiados; esta
identidad sólo delimita la identificación con una matriz unitaria finita.

Así queda articulado un conjunto de transportes concretos. La memoria
produce una carga y una acción parabólica; el marco de la elipse produce
la conexión variacional y su lagrangiano; las respuestas constitutivas
determinan los coeficientes de Maxwell; la graduación de canales construye
Clifford y permite factorizar el operador material; la reconstrucción de
Yang--Mills aporta una inmersión modal simpléctica que conserva exactamente
su brecha y su refinamiento. Cada identificación dispone de una aplicación,
una forma conservada y una ecuación entrelazada.
